\subsection{Проектирование системы}
\label{sec:design:system}

До начала реализации нужно определить границы системы: какие действия выполняет пользователь, что система делает сама по расписанию и с какими внешними сервисами обменивается данными. Эти границы описаны диаграммами вариантов использования и контекста. Поведение системы во время работы показано диаграммой последовательности и схемой алгоритма мониторинга. Диаграммы строились как рабочий материал для реализации: по ним определялось, какие данные читает и записывает каждая часть системы и в каком порядке.

\subsubsection{Диаграмма вариантов использования}

Диаграмма вариантов использования (\textit{Use Case Diagram}) описывает функциональные требования к системе с точки зрения тех, кто с ней взаимодействует. На ней показано, что система делает для пользователя, без деталей реализации. Диаграмма приведена на рисунке~\ref{fig:design:usecase}.

\begin{figure}[H]
	\centering
	\includegraphics[width=0.95\linewidth]{kufar/usecase}
	\caption{Диаграмма вариантов использования}
	\label{fig:design:usecase}
\end{figure}

На диаграмме выделены следующие действующие лица:
\begin{itemize}
	\item «Пользователь»~-- человек, который работает с ботом в \textit{Telegram}: добавляет, просматривает, изменяет и удаляет отслеживаемые товары, получает уведомления о выгодных предложениях и оценивает их, выгружает и загружает список товаров файлом, формирует отчёты, меняет язык интерфейса и привязывает электронную почту;
	\item «Администратор»~-- пользователь, идентификатор которого указан в конфигурации системы; наследует все возможности «Пользователя», получает повышенный лимит отслеживаемых товаров, отчёт о состоянии системы и служебные оповещения о сбоях;
	\item «\textit{Kufar API}»~-- внешняя система, поисковый интерфейс площадки. Из него система получает объявления при мониторинге и образцы объявлений при обновлении справочника фильтров;
	\item «Почтовый сервер»~-- внешняя система, через которую отправляются коды подтверждения, копии уведомлений и отчёты.
\end{itemize}

Названия вариантов использования заданы глаголами, поскольку вариант описывает действие, а не объект. Операции над отслеживаемым товаром образуют полный набор: «Добавить товар», «Просмотреть список товаров», «Просмотреть карточку товара», «Изменить товар» и «Удалить товар». Такой набор из пяти операций принято обозначать сокращением \textit{CRUDL} по первым буквам английских названий создания, чтения, изменения, удаления и получения списка.

Варианты «Опросить площадку» и «Обновить справочник» запускаются не человеком, а планировщиком системы по расписанию. Добавление товара включает выбор фильтров категории и проверку лимита, изменение товара включает тот же выбор фильтров, а импорт фильтров из ссылки \textit{Kufar.by} расширяет добавление как альтернативный способ задать критерии. Привязка электронной почты включает её подтверждение кодом, без которого письма на адрес не отправляются.

Модель вариантов использования разграничивает функции системы между двумя ролями. Пользователю доступна вся работа с отслеживаемыми товарами, уведомлениями, отчётами и почтой. Администратор получает эти же функции и, кроме них, средства контроля: отчёт о состоянии системы и оповещения о сбоях. Большинство вариантов использования заканчивается записью в базу данных, а получение уведомлений целиком построено на данных, накопленных предыдущими циклами опроса.

\subsubsection{Диаграмма контекста}

Диаграмма контекста первого уровня модели \textit{C4} показывает систему как единое целое вместе с людьми и внешними системами, с которыми происходит обмен данными. Внутреннее устройство на этом уровне не раскрывается [5]. Диаграмма приведена на рисунке~\ref{fig:design:context}.

\begin{figure}[H]
	\centering
	\includegraphics[width=0.6\linewidth]{kufar/context}
	\caption{Диаграмма контекста системы (\textit{C4}, уровень 1)}
	\label{fig:design:context}
\end{figure}

В центре диаграммы находится система отслеживания цен. Пользователь и администратор обращаются к ней не напрямую, а через \textit{Telegram}: команды и нажатия кнопок попадают в \textit{Telegram Bot API}, откуда система забирает их методом длинного опроса (\textit{long polling}). Тем же каналом система отправляет ответы, уведомления и служебные оповещения [6].

Площадка \textit{Kufar.by} выступает источником данных. Система обращается к её поисковому интерфейсу по \textit{HTTPS} и получает объявления в формате \textit{JSON}. Обращения ограничены: опрос выполняется раз в 3 минуты, и за минуту отправляется не более 30 запросов.

Все данные системы хранит СУБД \textit{PostgreSQL}. Система обращается к ней SQL-запросами. Почтовый сервер получает от системы письма с кодами подтверждения, копиями уведомлений и отчётами по протоколу \textit{SMTP}.

Граница системы проходит так, что ни одна внешняя система не обращается к ней сама по себе. Система опрашивает \textit{Telegram} и площадку по своей инициативе и в своём темпе. Это упрощает защиту: снаружи доступен только программный интерфейс для бота, закрытый секретным ключом.

\subsubsection{Диаграмма последовательности}

Диаграмма последовательности (\textit{Sequence Diagram}) показывает порядок обмена сообщениями между участниками во времени. Работа системы разделена на две независимые части, которые запускает планировщик: цикл опроса площадки и доставку уведомлений. Каждая часть показана отдельной диаграммой на рисунках~\ref{fig:design:sequence:scrape} и~\ref{fig:design:sequence:deliver}.

\begin{figure}[H]
	\centering
	\includegraphics[width=\linewidth]{kufar/sequence_scrape}
	\caption{Диаграмма последовательности цикла опроса площадки}
	\label{fig:design:sequence:scrape}
\end{figure}

Цикл опроса выполняется каждые 3 минуты. Задача опроса выбирает из базы данных активные отслеживаемые товары и для каждого отправляет площадке поисковый запрос с категорией, фильтрами и ограничением цены сверху, равным порогу товара. При отказе площадки с кодом 403 или 429 цикл прерывается сразу, не дожидаясь ответа по остальным товарам, и задача увеличивает паузу до следующего опроса.

При успешном ответе объявления передаются модулю отбора. Для каждого объявления модуль ищет в базе данных запись по паре «площадка~-- внешний идентификатор». Если записи нет, объявление добавляется, если есть, обновляется, и модуль получает прежнюю цену. Затем объявление проверяется по цене, словам-исключениям и названию. Совпадение создаётся, только если объявление опубликовано после добавления товара или его цена снизилась. Изменения всего цикла фиксируются одной транзакцией.

\begin{figure}[H]
	\centering
	\includegraphics[width=\linewidth]{kufar/sequence_deliver}
	\caption{Диаграмма последовательности доставки уведомлений}
	\label{fig:design:sequence:deliver}
\end{figure}

Доставка выполняется каждую минуту отдельной задачей. Задача выбирает недоставленные совпадения, у которых меньше пяти неудачных попыток и возраст не больше трёх суток, проверяет часовой и суточный лимиты и цену прежних отправок этого объявления. Если объявление уже приходило пользователю по такой же или меньшей цене, совпадение закрывается как дубликат. Иначе сообщение уходит через \textit{Telegram Bot API}. При успехе совпадение отмечается доставленным и в журнал уведомлений добавляется запись, при ошибке увеличивается счётчик попыток и сохраняется текст ошибки.

Опрос и доставка разделены намеренно. Медленный ответ \textit{Telegram} не задерживает опрос площадки, а совпадение, которое не удалось отправить, остаётся в базе данных и не теряется вместе с циклом. Очередью уведомлений в этой схеме служит обычная таблица совпадений, в которой недоставленные записи отличаются пустым временем доставки.

\subsubsection{Схемы алгоритмов мониторинга и доставки уведомлений}

Схемы алгоритмов детализируют логику, которая на диаграммах последовательности показана сообщениями. Схемы выполнены по ГОСТ~19.701–90: действие обозначено прямоугольником, условие~-- ромбом, начало и конец~-- прямоугольником со скруглёнными углами, ввод и вывод данных~-- параллелограммом. Символ «Начало» стоит вверху по центру, символ «Конец» на каждой схеме один и стоит внизу по центру. Линии потока проведены только по горизонтали и вертикали; стрелка поставлена там, где поток идёт не сверху вниз и не слева направо, то есть на возвратах к началу цикла и на боковых ветвях. Схема алгоритма цикла опроса приведена на рисунке~\ref{fig:design:algorithm:scrape}.

\begin{figure}[H]
	\centering
	\includegraphics{kufar/algorithm_scrape}
	\caption{Схема алгоритма цикла опроса площадки}
	\label{fig:design:algorithm:scrape}
\end{figure}

Алгоритм начинается с проверки паузы. Если на одном из предыдущих циклов площадка отказала в обслуживании, текущий цикл пропускается и счётчик пропускаемых циклов уменьшается на единицу. Пауза растёт при каждом отказе подряд: 1, 2, 4, 8 и не более 16 циклов. Первый успешный цикл после блокировки сбрасывает её. При интервале в 3 минуты самая длинная пауза составляет 48 минут.

Затем выбираются активные отслеживаемые товары, и для каждого формируется запрос к площадке. Товары обрабатываются строго по очереди, и нагрузка на площадку остаётся в пределах 30 запросов в минуту. Отказ площадки на любом товаре завершает весь цикл с оповещением администратора. Продолжать запросы в прежнем темпе после отказа означало бы превратить временное ограничение в длительную блокировку.

Обработка каждого полученного объявления вынесена в отдельную схему на рисунке~\ref{fig:design:algorithm:listing}. Объявление сначала сохраняется в базе данных, даже если потом будет отклонено: сохранённая цена нужна следующим циклам, потому что только сравнение с ней показывает снижение цены. Проверка релевантности состоит из трёх условий: цена ниже порога, в тексте нет слов-исключений, в заголовке или описании есть все слова названия товара. Прошедшее проверку объявление становится совпадением, только если оно опубликовано после добавления товара или подешевело. Объявления, которые были на площадке до добавления товара и с тех пор не подешевели, пользователь мог просмотреть сам при настройке поиска, а самые дешёвые из них получил при мгновенном поиске, поэтому о них система не сообщает.

\begin{figure}[H]
	\centering
	\includegraphics{kufar/algorithm_listing}
	\caption{Схема алгоритма обработки полученного объявления}
	\label{fig:design:algorithm:listing}
\end{figure}

Очередь доставки обрабатывает схема на рисунке~\ref{fig:design:algorithm:deliver}. Для каждого совпадения проверяются три условия по порядку: не исчерпан ли лимит уведомлений, не отправлялось ли объявление по такой же или меньшей цене и принял ли \textit{Telegram} сообщение. Первое условие оставляет совпадение в очереди до следующего часа или дня, второе закрывает его как дубликат, третье определяет, будет ли совпадение отмечено доставленным или получит ещё одну попытку. Изменения каждого запуска доставки фиксируются транзакцией, поэтому отметка о доставке и запись в журнале уведомлений сохраняются вместе.

\begin{figure}[H]
	\centering
	\includegraphics{kufar/algorithm_deliver}
	\caption{Схема алгоритма доставки уведомлений}
	\label{fig:design:algorithm:deliver}
\end{figure}
